\chapter{Introducción}
\label{cap:introduccion}

\section{Contexto y motivación}

Esta memoria describe el trabajo de la asignatura Procesos de Ingeniería del Software (PIS), curso 2026-2027, del Grado en Ingeniería Informática de la Escuela Superior de Ingeniería Informática de Albacete (Universidad de Castilla-La Mancha). El trabajo se plantea como si fuera un Trabajo Fin de Grado: un producto software real, desarrollado con un proceso de ingeniería (Scrum, control de versiones, integración y despliegue continuos y pruebas automáticas) y documentado con una memoria.

El enunciado del Sprint 1 fija cuatro hitos comunes a todos los proyectos. Son la base de cualquier aplicación SaaS: gestión de usuarios, autenticación, persistencia, roles, seguridad y despliegue. Sobre esa base, cada proyecto construye su propio producto. En este caso el producto es \textit{La Grada Colchonera}, un juego web para aficionados del Atlético de Madrid. Es un proyecto de aficionado, sin ninguna relación oficial con el Club Atlético de Madrid.

La motivación es doble. Por un lado, practicar el proceso completo de un TFG con un proyecto que se pueda usar de verdad: una URL pública, cuentas reales y datos guardados en una base de datos. Por otro, que el proceso se pueda revisar entero desde el repositorio: el backlog, las ramas, los pull requests, las pruebas y los despliegues dejan rastro, y la memoria se apoya en ese rastro.

\section{Objetivos}

\subsection{Objetivo principal}

Construir y desplegar una aplicación web pública en la que los aficionados del Atlético de Madrid se registran, entran (con email o con Google) y juegan a un juego propio, con su progreso guardado en el servidor, siguiendo un proceso de ingeniería del software que se pueda revisar de principio a fin.

\subsection{Objetivos específicos}

\begin{enumerate}
    \item \textbf{Producto.} Ofrecer la base SaaS que pide el enunciado (cuentas, sesión, roles, persistencia, correo, Google y manejo de errores) y, encima, el juego: una partida \textit{survivor} con un club que se guarda en el servidor, ranking, fichajes, entrenamiento y logros.
    \item \textbf{Proceso.} Aplicar Scrum con un backlog en DoneTonic, GitHub Flow con una rama y un pull request por tarea, integración continua con pruebas automáticas y despliegue continuo en la nube.
    \item \textbf{Calidad.} Usar una arquitectura en capas con responsabilidades separadas, probar los casos de error y no solo el camino feliz, aplicar seguridad desde el diseño (hash de contraseñas, secretos fuera del código y permisos comprobados en el servidor) y hacer una interfaz accesible y adaptada al móvil.
    \item \textbf{Documentación.} Que todo el trabajo, incluida la gestión, se pueda revisar desde el repositorio: README, carpeta \path{docs/} y esta memoria.
\end{enumerate}

\section{Alcance y cambio de enfoque}

El proyecto se ha desarrollado en dos sprints:

\begin{itemize}
    \item \textbf{Sprint 1} (29/09/2026 -- 27/10/2026): los cuatro hitos del enunciado, con 19 PBIs (P1--P19) y 32 tareas. Los Hitos 1--3 (repositorio, CI, despliegue continuo, lógica de usuarios con pruebas, API REST, frontend separado en interfaz y cliente, registro e inicio de sesión, sesión, rutas protegidas, PostgreSQL, secretos, roles, panel de administración y registro de actividad) se entregaron el 12/10/2026. El Hito 4 (confirmación de la cuenta por correo, inicio de sesión con Google, manejo de errores de principio a fin y política de privacidad) se completó el 23/10/2026. Al cierre del sprint se añadió el despliegue en Google Cloud Run.
    \item \textbf{Sprint 2} (23/10/2026 -- 29/10/2026): la lógica del juego, con 29 PBIs (P20--P48) y 50 tareas: el rediseño de la interfaz y el juego, construido en seis fases.
\end{itemize}

\subsection{Por qué un juego y no una red social}

La idea inicial era una red social para aficionados del Atlético. Cerrado el Sprint 1, el profesor dejó libertad para la lógica y el diseño del producto (y fijó después que el Sprint 2 sería esa lógica), y se decidió hacer un juego por tres motivos:

\begin{enumerate}
    \item Es más original que otro clon de red social.
    \item Obliga a resolver problemas técnicos más ricos: simulación determinista, rendimiento al dibujar en \textit{canvas}, generación procedural del mapa, sonido sintetizado y validación de las partidas en el servidor para evitar trampas.
    \item Reutiliza entera la base del Sprint 1. Un juego con cuentas, progreso en el servidor, ranking y roles necesita exactamente lo que piden los hitos.
\end{enumerate}

Quedan fuera del alcance, por ahora, el multijugador en tiempo real, los amigos y las partidas en grupo, y las aplicaciones nativas (Sección~\ref{sec:trabajo-futuro}).

\section{Estructura de la memoria}

\begin{itemize}
    \item El Capítulo~\ref{cap:propuesta} presenta la propuesta, la metodología (Scrum en DoneTonic y GitHub Flow) y las herramientas.
    \item El Capítulo~\ref{cap:descripcion} describe la aplicación: requisitos, arquitectura, modelo de datos, API, los flujos principales con diagramas de secuencia, la seguridad y el uso de la aplicación con capturas.
    \item El Capítulo~\ref{cap:juego} explica el juego del Sprint 2: concepto, cómo se juega, sistemas, arquitectura, servidor y anti-trampas, logros y evolución por fases.
    \item El Capítulo~\ref{cap:evaluacion} recoge la evaluación: pruebas, integración y despliegue continuos, registro de actividad y errores, revisión frente al enunciado, prueba técnica, comprobación en producción y métricas.
    \item El Capítulo~\ref{cap:conclusiones} repasa los objetivos, las lecciones aprendidas y el trabajo futuro.
    \item El Anexo~\ref{anx:desarrollo} cuenta el desarrollo de cada sprint; el Anexo~\ref{anx:trazabilidad}, la trazabilidad requisito -- PBI -- tarea -- pull request; el Anexo~\ref{anx:material}, los enlaces al material; y el Anexo~\ref{anx:ia}, el uso de la inteligencia artificial.
\end{itemize}
